\documentclass[11pt]{article}

\usepackage[T1]{fontenc}
\usepackage[margin=1in]{geometry}
\usepackage{hyperref}
\usepackage{booktabs}
\usepackage{amsmath,amssymb}

\title{HDG Diffusion-Reaction Flux Postprocessing}
\author{}
\date{}

\begin{document}
\maketitle

\section{Context}

The implementation and verification path is exercised by the preset
\begin{quote}
\small
\begin{verbatim}
scripts/diffusion_reaction/run_cases.py trigonometric_poisson_cg_gamg \
    --verbosity 2 --plot
\end{verbatim}
\end{quote}
where the scalar postprocessing was very effective, but the flux postprocessing
was only modestly better than the raw HDG flux:
\[
\|q_h - q\|_{L^2} = 4.1124 \times 10^{-9}, \qquad
\|q_h^\star - q\|_{L^2} = 2.1696 \times 10^{-9}.
\]

The relevant files were:
\begin{itemize}
    \item \verb|scripts/diffusion_reaction/run_cases.py| for preset execution and error reporting,
    \item \verb|hdgfem/solvers/diffusion_reaction.py| for the HDG solver and postprocessing wrapper,
    \item \verb|hdgfem/mixed/postprocess/numba_kernels.py| (current location) for the fused Numba local postprocessing kernels,
    \item \verb|tests/test_diffusion_reaction_solver.py| for solver and postprocessing regression tests.
\end{itemize}

\section{Diagnosis}

The runner's vector $L^2$ error calculation was not the source of the issue.
For a postprocessed flux, it evaluates the exact flux on the postprocessed
degree $p+1$ quadrature space, so the comparison is made at the correct
quadrature order.

The implemented flux postprocessor was mathematically consistent with its
documented constraints. It built a degree $p+1$ vector field whose:
\begin{itemize}
    \item normal moments match the HDG numerical flux
    \[
    q_h \cdot n + \tau (u_h - \widehat{u}_h),
    \]
    \item interior component moments match the raw HDG flux against
    $[P_{p-1}]^2$,
    \item remaining unconstrained high-order modes are chosen by a minimum
    mass-norm correction from an $L^2$ projection of raw $q_h$.
\end{itemize}

That last point explained the limited improvement. The constrained projection
was preserving conservation constraints, but its free high-order content was
still tied to raw $q_h$ rather than to the much more accurate postprocessed
scalar field $u_h^\star$.

\section{Postprocessing Formulas}

The mixed HDG variables use the conservative flux convention
\[
q = -\kappa \nabla u.
\]
On each element $K$, the computed local fields are
\[
u_h \in P_p(K), \qquad q_h \in [P_p(K)]^2,
\]
with trace unknown $\widehat{u}_h$ on the element faces. The local
postprocessing constructs degree $p+1$ fields
\[
u_h^\star \in P_{p+1}(K), \qquad
q_h^\star \in [P_{p+1}(K)]^2.
\]

\subsection*{Primal Postprocessing}

The scalar postprocessed field is the standard local Neumann recovery: find
$u_h^\star \in P_{p+1}(K)$ such that
\begin{align}
(\nabla u_h^\star, \nabla w)_K
    &= -(\kappa^{-1} q_h, \nabla w)_K
    &&\forall w \in P_{p+1}(K), \label{eq:primal-post-gradient}\\
(u_h^\star, 1)_K
    &= (u_h, 1)_K. \label{eq:primal-post-mean}
\end{align}
Equation \eqref{eq:primal-post-gradient} recovers a scalar field whose
gradient is consistent with the mixed flux relation, while
\eqref{eq:primal-post-mean} fixes the constant mode.

\subsection*{Flux Postprocessing}

The flux postprocessor is a constrained local projection. Let $q_0$ be the
reference flux used for the minimum-distance projection. The postprocessed flux
is defined as
\[
q_h^\star
    = \arg\min_{v \in [P_{p+1}(K)]^2}
        \frac{1}{2}\|v - q_0\|_{0,K}^2
\]
subject to the HDG conservation constraints
\begin{align}
\langle q_h^\star \cdot n, \mu \rangle_F
    &=
    \langle q_h \cdot n + \tau (u_h - \widehat{u}_h), \mu \rangle_F
    &&\forall \mu \in P_{p+1}(F), \quad F \subset \partial K,
    \label{eq:flux-post-normal}\\
(q_h^\star, r)_K
    &= (q_h, r)_K
    &&\forall r \in [P_{p-1}(K)]^2.
    \label{eq:flux-post-interior}
\end{align}
The first constraint matches the numerical normal flux on every face. The
second preserves the low-order interior flux moments.

In matrix form, with vector mass matrix $A$, constraint matrix $C$, and
constraint right-hand side $b$, the constrained minimizer satisfies
\begin{align}
q_h^\star &= q_0 + A^{-1} C^T \lambda, \\
(C A^{-1} C^T)\lambda &= b - C q_0.
\end{align}
The existing implementation already factored the element-local Schur
complement $C A^{-1} C^T$. The improvement made in this chat changes only
the reference $q_0$ in the identity-diffusion \verb|both| path:
\[
q_0 =
\begin{cases}
\Pi_{p+1} q_h,
    & \text{flux-only mode or tensor-diffusion fallback}, \\
\Pi_{p+1}(-\nabla u_h^\star),
    & \text{identity diffusion with primal and flux postprocessing}.
\end{cases}
\]
Here $\Pi_{p+1}$ denotes the elementwise $L^2$ projection into
$[P_{p+1}(K)]^2$. Thus the new path leaves
\eqref{eq:flux-post-normal}--\eqref{eq:flux-post-interior} unchanged, but it
selects the unconstrained high-order part of $q_h^\star$ from the more
accurate scalar recovery.

subsection*{Selectable Raviart--Thomas Projection}

The alternative erb|flux_postprocess_space="RT_projection"| defines
$q_h^{mathrm{RT}}$ as the unique member of
[
    [P_p(K)]^2 + x P_p(K)
]
satisfying
egin{align}
langle q_h^{mathrm{RT}}cdot n,mu
angle_F
  &= langle q_hcdot n+	au(u_h-widehat u_h),mu
angle_F
  &&orall muin P_p(F),quad Fsubsetpartial K,\
(q_h^{mathrm{RT}},r)_K
  &= (q_h,r)_K
  &&orall rin[P_{p-1}(K)]^2,quad pgeq1.
end{align}
For $p=0$, only the face-normal moments are present. These are the unisolvent
Raviart--Thomas degrees of freedom. Because the HDG numerical normal flux is
single-valued on interior faces, the elementwise reconstruction belongs to
$H(operatorname{div},Omega)$. Its two polynomial components are stored
exactly in the existing degree-$p+1$ DG vector space; that storage choice does
not enlarge the reconstructed RT space.

The RT moment system can be solved with the host Numba kernel or the batched
CuPy implementation. The default erb|l2_closest| method remains the
full-space constrained recovery above. Compatibility spellings
erb|full-p-plus-1| and erb|rt-p| map to those two canonical choices.

\section{Implementation}

For identity diffusion and \verb|hdg_postprocess="both"|, the flux
postprocessor now keeps the same HDG constraints but changes the
minimum-distance reference from raw $q_h$ to the constitutive flux
\[
q_0 = -\nabla u_h^\star.
\]

This affects only the identity-diffusion \verb|both| path. Flux-only
postprocessing and tensor-diffusion postprocessing still use the original
raw-$q_h$ reference path.

The implementation changes were:
\begin{itemize}
    \item Added cached reference projection matrices
    \[
    M^{-1}(\phi_i, \partial_r \phi_j), \qquad
    M^{-1}(\phi_i, \partial_s \phi_j)
    \]
    in \verb|_HDGPostprocessCache|, now in \verb|hdgfem.mixed.postprocess.flux|.
    \item Added the primal-reference H(div) flux kernel in
    the Numba postprocessing kernels, now \verb|hdgfem.mixed.postprocess.numba_kernels|:
    \begin{quote}
    \small
    \verb|solve_hdiv_flux_primal_reference_min_distance_postprocess_kernel|
    \end{quote}
    \item Factored the shared constraint projection into
    \begin{quote}
    \small
    \verb|_apply_hdiv_flux_min_distance_constraints|
    \end{quote}
    so both the old and new
    flux references use exactly the same normal-flux and interior-moment
    constraints.
    \item Routed identity-diffusion \verb|both| postprocessing through the new
    primal-reference kernel.
\end{itemize}

\section{Verification}

The focused regression command is:
\begin{quote}
\verb|.venv/bin/python -m pytest tests/test_diffusion_reaction_solver.py -q|
\end{quote}
The tests verify that the postprocessed flux improves on raw $q_h$ for a smooth
identity-diffusion problem while preserving the H(div)-style moment
constraints. They also cover the tensor-diffusion fallback and require
degree-$p+1$ primal and flux outputs.

\section{Observed Results}

On a small structured p=3 sequence, the improved flux was consistently lower:
\begin{center}
\begin{tabular}{rrrr}
\toprule
$n_x$ & elements & raw flux error & post flux error \\
\midrule
4  & 32  & $3.052875 \times 10^{-3}$ & $1.026032 \times 10^{-3}$ \\
8  & 128 & $2.279693 \times 10^{-4}$ & $6.146711 \times 10^{-5}$ \\
12 & 288 & $4.635780 \times 10^{-5}$ & $1.212750 \times 10^{-5}$ \\
\bottomrule
\end{tabular}
\end{center}

On the reported PETSc/GAMG preset, the postprocessed flux error improved from
the original
\[
2.1696 \times 10^{-9}
\]
to
\[
1.1386 \times 10^{-9}.
\]
The raw flux error remained
\[
4.1124 \times 10^{-9},
\]
and the postprocessed primal error remained
\[
2.9877 \times 10^{-11}.
\]

Additional PETSc/GAMG runs after the primal-reference flux change showed the
same pattern on nearby mesh sizes. The 72,969-triangle summary appeared twice
in the pasted logs and is listed once here.
\begin{center}
\small
\setlength{\tabcolsep}{4pt}
\begin{tabular}{rrrrrrr}
\toprule
triangles & trace dofs & $h^{p+1}$ & raw flux & post flux & reduction & post primal \\
\midrule
28,569 & 301,350 & $1.5872{\times}10^{-7}$ & $3.0769{\times}10^{-8}$ &
$8.2495{\times}10^{-9}$ & $3.73{\times}$ & $1.0951{\times}10^{-10}$ \\
72,969 & 768,376 & $6.8735{\times}10^{-9}$ & $1.1368{\times}10^{-9}$ &
$3.0657{\times}10^{-10}$ & $3.71{\times}$ & $1.4182{\times}10^{-11}$ \\
\bottomrule
\end{tabular}
\end{center}
Both runs used $p=6$, identity diffusion, $\tau=1$, eliminated boundary
trace degrees of freedom, the Numba assembly backend, and PETSc CG/GAMG with
10 levels. Their reported postprocessing costs were $0.6$ seconds and
$1.6$ seconds, respectively, about $5$--$6\%$ of total runtime.

\end{document}
